\documentclass[11pt]{amsart}
\usepackage[T1]{fontenc}
\usepackage[utf8]{inputenc}
\usepackage{amssymb}
\usepackage[margin=1.1in]{geometry}
\usepackage[hidelinks]{hyperref}
\newtheorem{theorem}{Theorem}
\newtheorem{proposition}[theorem]{Proposition}
\newtheorem{lemma}[theorem]{Lemma}
\newtheorem{corollary}[theorem]{Corollary}
\theoremstyle{definition}
\newtheorem{definition}[theorem]{Definition}
\newtheorem{axiom}[theorem]{Axiom}
\theoremstyle{remark}
\newtheorem{remark}[theorem]{Remark}
\newcommand{\Om}{\Omega}
\newcommand{\Map}{\mathrm{Map}_*}
\newcommand{\Z}{\mathbb{Z}}
\newcommand{\code}[1]{\texttt{#1}}
\newcommand{\Hom}{\mathrm{Hom}}
\newcommand{\hc}[2]{[#1,#2]_*}
\newcommand{\Fib}{\mathrm{Fib}}
\newcommand{\tr}[1]{\|#1\|}
\newcommand{\cls}[1]{[#1]}
\title{$\pi_3(S^2)=\Z$ by a Postnikov universal property:\\ the main ideas of a formal proof}
\date{28 October 2026}
\begin{document}
\maketitle

\noindent\textbf{Attribution.} David W\"arn developed the
universal-property strategy used here (one identifies a candidate Postnikov stage by a
\emph{marked universal property} and compares it with the truncated sphere on homotopy
classes), together with the delooping constructions (Eilenberg--MacLane spaces as
iterated torsors, and the criterion for a map of loop spaces to be the looping of a
map~\cite{DW1}) and the cup-product construction. Our contribution was to develop the
argument in detail and formalise it in Unit; the formal proof is the development
\code{sag/pi3-s2-postnikov}, and the detailed notes are listed in \S\ref{s:formal}.
Owen Milner formulated the S\'{\i}nh invariant in homotopy type theory.

\smallskip\noindent\textbf{The result.} In univalent type theory with replacement,
and assuming only the universal property of $S^2$ (Axiom~\ref{a:s2}),
\[\pi_3(S^2)\cong\Z\quad\text{as groups}\]
(\code{Presentation.Final.pi3-S2}, Theorem~\ref{t:final}). No additional
higher-inductive construction is assumed beyond the stated universal property of the
sphere; no pushout, no Hopf fibration, no smash product and no spectral sequence is
used. The homotopy-theoretic input is the cohomology of $K(\Z,2)$ in degrees 3 and 4,
proved from the circle and the cup square. The same argument one level up gives
$\pi_4(S^3)=\Z/2$ (\code{doc/pi4-s3.pdf}).

\section{Setting}

The foundations are those of \code{sag/em}: univalence and Rijke's replacement
(images and their surjectivity) as parameters; function extensionality,
propositional truncation, set quotients and the finite truncations $\tr X_m$ are
derived. Eilenberg--MacLane spaces $K(A,n)$ of an abelian group $A$ are built as in
\cite{DW1}: $K(A,1)$ is the type of $A$-torsors and $K(A,n{+}1)$ is the type of
$K(A,n)$-torsors, i.e.\ of merely inhabited types with a homogeneous structure. We
write $e_n:\Om K(A,n{+}1)\simeq_*K(A,n)$, $K(h,n)$ for the map induced by a
homomorphism $h$, and $F;G$ for composition of pointed maps (first $F$).

For pointed $X,Y$, $\hc XY$ is the \emph{set} of pointed homotopy classes, the set
truncation of $\Map(X,Y)$; classes compose and form a category. Two classes are equal
when their representatives are merely equal; no representative is ever chosen. The
third homotopy group is $\pi_3X=\pi_1(\Om^2X)$, functorial on maps and on classes.

\begin{axiom}[$S^2$]\label{a:s2}
A type $S^2$ with a point $\mathsf{base}$ and $\mathsf{surf}:\Om^2(S^2,\mathsf{base})$
such that for every $X$, $f\mapsto(f\,\mathsf{base},\mathsf{ap}^2_f\,\mathsf{surf})$ is
an equivalence $(S^2\to X)\simeq(x:X)\times\Om^2(X,x)$.
\end{axiom}

The axiom is a hypothesis of the module \code{SphereAxiom.Sphere} (its parameters
\code{S2}, \code{base}, \code{surf}, \code{ump}), not of the whole development, so only
the results that take these parameters depend on it. From it one proves that the
pointed evaluation $\Map(S^2,X)\to\Om^2X$, $F\mapsto(\Om^2F)(\mathsf{surf})$, is an
equivalence, and that it is natural: for $G:X\to_*Y$, the evaluation of $F;G$ is
$\Om^2G$ applied to the evaluation of $F$.

\section{The strategy: a marked universal property}

The proof is a Yoneda argument on classes.

Let $T=\tr{S^2}_3$ with unit $\eta:S^2\to_*T$. From the axiom, $S^2$ is simply
connected (into a 1-type $Y$ evaluation at the base point is an equivalence, since
$\Om^2Y$ is contractible, so $\tr{S^2}_1$ is contractible and $S^2$ is the fibre of its
unit). Hence $T$ is a simply connected 3-type, $\eta^*:\Map(T,Y)\to\Map(S^2,Y)$ is an
equivalence of pointed mapping types for every 3-type $Y$, hence a bijection
$\hc TY\to\hc{S^2}Y$ on homotopy classes, and $\pi_3\eta$ is an isomorphism (the fibres
of $\Om^k\eta$ are $(3-k)$-connected).

\begin{proposition}[two instances suffice]\label{p:two}
Let $P$ be a 3-type and $s\in\hc{S^2}P$ such that $s^*:\hc PY\to\hc{S^2}Y$ is a
bijection for $Y=P$ and for $Y=T$. Then $\pi_3S^2\cong\pi_3T\cong\pi_3P$.
\end{proposition}
\begin{proof}
Let $u\in\hc TP$ with $\cls\eta;u=s$ ($\eta^*$ at $P$) and $t\in\hc PT$ with
$s;t=\cls\eta$ ($s^*$ at $T$). Then $s;(t;u)=\cls\eta;u=s=s;\mathrm{id}$, so
$t;u=\mathrm{id}_P$; and $\cls\eta;(u;t)=s;t=\cls\eta$, so $u;t=\mathrm{id}_T$. Hence
$\pi_3u$ and $\pi_3t$ are inverse isomorphisms.
\end{proof}

So the proof consists in constructing a 3-type $P$ with $\pi_3P=\Z$ and a class $s$
with this universal property. The candidate is $P=\Fib\kappa$, the fibre of the cup
square $\kappa:K(\Z,2)\to_*K(\Z,4)$; this is the Postnikov stage with $k$-invariant
$\kappa$, and $\pi_3P\cong\Z$ as groups since the fibre of $P\to K(\Z,2)$ is
$\Om K(\Z,4)\simeq K(\Z,3)$.

\section{The cup square and the cohomology of $K(\Z,2)$}

\noindent\textbf{The circle.} $K(\Z,1)$ has the universal property of the circle:
maps $K(\Z,1)\to Y$ correspond to a point and a loop of $Y$, following the
construction of the circle as $\Z$-torsors of Bezem, Buchholtz, Grayson and
Shulman~\cite{BBGS} (\code{doc/circle-torsors.pdf}).

\noindent\textbf{The cup square.} Following W\"arn's construction, the circle property
and delooping give a bipointed map $c_2:K(\Z,2)\to_*\Map(K(\Z,2),K(\Z,4))$, normalised
so that looping it and evaluating at the generator of $\pi_1K(\Z,1)$ gives back the
identity of $K(\Z,2)$; the cup square is its diagonal, $\kappa(x)=c_2(x)(x)$
(\code{doc/cup-square.pdf}).

\noindent\textbf{Cohomology of $K(\Z,2)$.} For every abelian group $B$:
\begin{itemize}
\item[(H3)] $\hc{K(\Z,2)}{K(B,3)}=0$. Looping has propositional fibres
  on $\Map(K(\Z,2),K(B,3))$ \cite{DW1}, and its target is, up to $e_1,e_2$, the type
  $\Map(K(\Z,1),K(B,2))\simeq\Om K(B,2)\simeq K(B,1)$ (circle property), which is
  connected (\code{doc/circle-torsors.pdf});
\item[(H4)] $\Phi_B:\Hom(\Z,B)\to\hc{K(\Z,2)}{K(B,4)}$,
  $h\mapsto\cls{\kappa;K(h,4)}$, is a bijection (\code{doc/h4.pdf}).
\end{itemize}
(H4) is the composite of two bijections. (A) \emph{Every class is a diagonal}: for
$Y=K(B,4)$, $c\mapsto\cls{\Delta c}$ is a bijection from bipointed maps
$K(\Z,2)\to_*\Map(K(\Z,2),Y)$ to $\hc{K(\Z,2)}Y$; the proof embeds classes into a set
of ``multiplicativity defects'' by W\"arn's Yoneda step over the path fibration
$K(\Z,1)\to *\to K(\Z,2)$ (\cite{DW1}, Lemma~9), and computes the defect of a diagonal
by one path induction. (B) \emph{Bilinear classification}: $g\mapsto c_2;K(g,4)_*$ is
a bijection from $\Hom(\Z,B)$ to bipointed maps, from deloopings, the circle property
and the normalisation of $c_2$. Since $\Delta(c_2;K(g,4)_*)=\kappa;K(g,4)$ by
definition, $\Phi_B$ is the composite. No coefficient is ever recovered from a cross
term $c(x)(y)+c(y)(x)$, where a factor $2$ would appear.

\section{The mapping theorem for $P$}

A \emph{Postnikov triple} is $(A,B,k)$ with $k:K(A,2)\to_*K(B,4)$; its fibre $\Fib k$ is
a simply connected 3-type, with projection $q:\Fib k\to_*K(A,2)$.

\begin{theorem}[lifting]\label{t:lift}
For a pointed $Q$, suppose (E) every $F:Q\to_*K(A,2)$ has merely $F;k=0$, and (U) every
$d:Q\to_*\Om K(B,4)$ is merely $0$. Then $q_*:\hc Q{\Fib k}\to\hc Q{K(A,2)}$ is a
bijection.
\end{theorem}

Lifts exist for maps (by (E)), and two lifts with merely equal projections are merely
equal (by (U), and pointing irrelevance since $\Om\Fib k$ is connected); taking classes
is never assumed to preserve fibres (\code{doc/mapping-fib.pdf}). The lemma is applied
to two sources. For $Q=S^2$ both conditions are immediate: pointed maps out of $S^2$
are 2-loops, and $\Om^2K(B,4)$, $\Om^3K(B,4)$ are connected.

\noindent\textbf{(E) for $P$.} Let $p:P\to_*K(\Z,2)$ be the projection. Its fibres are
2-connected, so a map $F:P\to_*K(A,2)$ factors merely as $F=p;G$ with
$G:K(\Z,2)\to_*K(A,2)$. By (H4), $\cls{G;k}=\Phi_B(g)=\cls{\kappa;K(g,4)}$ for some $g$,
so merely $G;k=\kappa;K(g,4)$, and
\[F;k=p;(G;k)=(p;\kappa);K(g,4)=0;K(g,4)=0,\]
merely, since $p;\kappa$ is null by the definition of $P$: every obstruction to lifting
out of $P$ is a multiple of $\kappa$, and $\kappa$ is killed by $P$.

\noindent\textbf{(U) for $P$: $\hc P{K(B,3)}=0$.} This is the exactness of the Serre
sequence $0=H^3(K(\Z,2))\to H^3(P)\to H^3(K(\Z,3))\xrightarrow{\ \Phi_B\ }H^4(K(\Z,2))$,
proved for this fibration only, with explicit null-homotopies
(\code{doc/fibre-h3.pdf}). A map $F:P\to_*K(B,3)$ restricts on the fibre $\Om K(\Z,4)$ over
the base point to a map that deloops to some $G:K(\Z,4)\to_*K(B,4)$, and $G=K(g,4)$;
$F$ gives an explicit null-homotopy of $\kappa;G$ (constancy on the 2-connected fibres
is a proposition and spreads along $K(\Z,2)$), so $g=0$ by injectivity of $\Phi_B$ and
the restriction is null; then $F$ factors through
$K(\Z,2)$, and (H3) kills it. The connecting map is identified with $\Phi_B$ by
construction; no separate connecting map is defined.

\noindent\textbf{The marked class.} From the axiom and the circle property one builds
$\sigma:S^2\to_*K(\Z,2)$ with $(\Om^2\sigma)(\mathsf{surf})=(\Om e_1^{-1})(g)$, $g$ the
generator of $\pi_1K(\Z,1)$. Then $\sigma^*:\Map(K(\Z,2),Y)\to\Map(S^2,Y)$ is an
equivalence for every 2-type $Y$: up to the evaluation at $\mathsf{surf}$ it is the
composite of looping (an equivalence, since $K(\Z,2)$ is simply connected and $Y$ a
2-type), precomposition with $e_1^{-1}$, and the circle property. So $K(\Z,2)$ is the
2-truncation of $S^2$, pointed. Let
\[s=p_*^{-1}\cls\sigma\in\hc{S^2}P\]
(Theorem~\ref{t:lift} at $Q=S^2$ for the triple $(\Z,\Z,\kappa)$); $s$ is a class,
obtained without choice. Then $s^*$ is a bijection into every $K(A,2)$ (since
$s^*\circ p^*=\cls\sigma^*$, and $p^*$ is a bijection into 2-types), and the commuting
square of $q_*$ for the two sources gives:

\begin{theorem}[mapping theorem]\label{t:main}
For every Postnikov triple, $s^*:\hc P{\Fib k}\to\hc{S^2}{\Fib k}$ is a bijection; hence
so is $s^*:\hc PY\to\hc{S^2}Y$ for every $Y$ merely pointed equivalent to some
$\Fib k$.
\end{theorem}

\section{A Postnikov presentation of $T$}

\begin{theorem}\label{t:pres}
$T=\tr{S^2}_3$ is merely pointed equivalent to $\Fib k$ for a Postnikov triple
$(\Z,B,k)$.
\end{theorem}
\begin{proof}[Idea]
(1) $V=\tr T_2$ is pointed equivalent to $K(\Z,2)$: both have the universal property of
the 2-truncation of $S^2$ for maps into 2-types ($\eta$ and the unit of $V$ on one side,
$\sigma$ on the other). (2) The fibre $E_0$ of $u:T\to V$ over the base point is a
2-connected 3-type, hence $K(B,3)$ with $B=\Om^3E_0$: a 2-connected 3-type is
recognised by delooping its loop space twice. (3) $K(B,4)$ \emph{is} the type of
$K(B,3)$-torsors, so a map $V\to K(B,4)$ whose fibres are the fibres $E_z$ of $u$ is the
same as a homogeneous structure on every $E_z$. (4) The structures $K(B,3)\simeq_*E_z$
form a family of \emph{sets} over the simply connected $T$, so the one given at the base
point by (2) extends to all of $T$; simple connectivity removes the twisting. This
gives $k':V\to K(B,4)$ with $\Fib k'\simeq T$; take $k=\psi;k'$ with
$\psi:K(\Z,2)\simeq_*V$. The pointing is only mere, since $T$ is connected.
\end{proof}

The group $B$ is not assumed to be $\Z$: it is $\Om^3E_0$, and $B\cong\Z$ is a
consequence of the final comparison. Since the goal ``$s^*$ is a bijection'' is a
proposition, the mere presentation suffices.

\section{Conclusion}

\begin{theorem}\label{t:final}
Under Axiom~\ref{a:s2}, $\pi_3(S^2)\cong\Z$ as groups.
\end{theorem}
\begin{proof}
Theorem~\ref{t:main} applies to $Y=P=\Fib\kappa$ directly and to $Y=T$ through
Theorem~\ref{t:pres}. Proposition~\ref{p:two} gives
\[\pi_3S^2\cong\pi_3T\cong\pi_3P\cong\Z.\qedhere\]
\end{proof}

\section{The formalisation}\label{s:formal}

The proof is checked in Unit, in \code{sag/pi3-s2-postnikov} (\code{./segc
sag/pi3-s2-postnikov all}; a fresh check of the whole development, including
$\pi_4(S^3)$, has 4038 declarations, 0 failed, about 2\,s). The final declaration is
\begin{center}
\code{Presentation.Final[S2, base, surf, ump].pi3-S2 : GIso[pi3 S2 base, ab-grp Zab]}
\end{center}
whose only parameters beyond the foundations (univalence, images and their surjectivity)
are the $S^2$ axiom; the 3-truncation is constructed inside.

\smallskip
\begin{tabular}{@{}ll@{}}
\code{SphereAxiom} & the axiom\\
\code{HomotopyClasses}, \code{HomotopyGroups} & classes; $\pi_3$\\
\code{Truncations}, \code{SphereTruncation} & finite truncations; $T=\tr{S^2}_3$\\
\code{SphereUniversal} & Proposition~\ref{p:two}\\
\code{PostnikovRealization} & $\Fib k$ is a simply connected 3-type, $\pi_3\cong B$\\
\code{CircleProperty}, \code{H3Proof} & the circle; (H3)\\
\code{CupSquare}, \code{H4Diagonal}, \code{H4Bilinear}, \code{H4Proof} & the cup square; (A), (B), (H4)\\
\code{FibreH3} & $\hc P{K(B,3)}=0$\\
\code{MapFib} & Theorem~\ref{t:lift}; $\sigma$; $s$; Theorem~\ref{t:main}\\
\code{Presentation} & Theorem~\ref{t:pres}; Theorem~\ref{t:final}
\end{tabular}

\smallskip
For every declaration the checker reports which definitions it had to unfold; the
\code{.opens} files of the development record, declaration by declaration, the
definitions allowed to be unfolded, and the check against them passes for every
declaration. These specifications were generated from the checker's report and reviewed
by kind (which kinds of unfolding each text step justifies); the textual audit, entry by
entry against the notes, is not yet complete. The circle property, (H3), the cup square,
(H4), the fibre lemma, the mapping theorem for $\Fib\kappa$ and the presentation of $T$
were written from an informal note first; the earliest modules (the sphere axiom,
homotopy classes and groups, the realisation of Postnikov triples, finite truncations)
were written before their notes. Detailed notes: \code{circle-torsors.pdf},
\code{cup-square.pdf}, \code{h4.pdf}, \code{fibre-h3.pdf}, \code{mapping-fib.pdf},
\code{postnikov-presentation.pdf}, \code{finite-truncation.pdf}, and
\code{pi3-s2-postnikov.pdf} (the plan and its status).

\begin{thebibliography}{9}
\bibitem{BBGS} M.~Bezem, U.~Buchholtz, D.~R.~Grayson, M.~Shulman. \emph{Construction of
  the circle in UniMath}. arXiv:1910.01856, 2019.
\bibitem{DW1} D.~W\"arn. \emph{Eilenberg--MacLane spaces and stabilisation in homotopy
  type theory}. arXiv:2301.03685, 2023.
\end{thebibliography}
\end{document}
